\documentclass[
    12pt, % Schriftgröße
    DIV10,
    ngerman, % für Umlaute, Silbentrennung etc.
    a4paper, % Papierformat
    oneside, % einseitiges Dokument
    titlepage, % es wird eine Titelseite verwendet
    parskip=half, % Abstand zwischen Absätzen (halbe Zeile)
    headings=normal, % Größe der Überschriften verkleinern
    listof=totoc, % Verzeichnisse im Inhaltsverzeichnis aufführen
    bibliography=totoc, % Literaturverzeichnis im Inhaltsverzeichnis aufführen
    index=totoc, % Index im Inhaltsverzeichnis aufführen
    captions=tableheading, % Beschriftung von Tabellen unterhalb ausgeben
    final % Status des Dokuments (final/draft)
]{scrartcl}

% Pakete

% Schrift ----------------------------------------------------------------------
\usepackage{lmodern} % bessere Fonts

% Paket für Kopfzeilen und Fußzeilen
\usepackage[
    automark, % Kapitelangaben in Kopfzeile automatisch erstellen
    headsepline, % Trennlinie unter Kopfzeile
    ilines % Trennlinie linksbündig ausrichten
]{scrpage2}

% Anpassung an Landessprache
\usepackage[ngerman]{babel}

% Umlaute ----------------------------------------------------------------------
%   Umlaute/Sonderzeichen wie äöüß direkt im Quelltext verwenden (CodePage).
%   Erlaubt automatische Trennung von Worten mit Umlauten.
% ------------------------------------------------------------------------------
\usepackage[utf8]{inputenc}
\usepackage[T1]{fontenc}
\usepackage{textcomp} % Euro-Zeichen etc.

\usepackage[babel,german=quotes]{csquotes}

% Bei Änderungen müssen die *.aux und *.bbl Dateien manuell gelöscht werden
% Sonst kommt es zu einem Fehler bei der Erstellung
%\usepackage[round, sort, comma, numbers]{natbib}
%\usepackage[round, numbers]{natbib}
%\bibliographystyle{abbrvdin}
%

\usepackage{scrhack}

\usepackage[citestyle=alphabetic,bibstyle=alphabetic,backend=biber]{biblatex}
\bibliography{Bibliographie}
\DefineBibliographyStrings{german}{
  bibliography = {Literaturverzeichnis},
}
\setlength\bibitemsep{8pt}

\usepackage{xpatch}
\xpretobibmacro{author}{\mkbibbold\bgroup}{}{}
\xapptobibmacro{author}{\egroup}{}{}
\xpretobibmacro{editor}{\mkbibbold\bgroup}{}{}
\xapptobibmacro{editor}{\egroup}{}{}
\xpretobibmacro{editor+others}{\mkbibbold\bgroup}{}{}
\xapptobibmacro{editor+others}{\egroup}{}{}
\xpretobibmacro{bbx:editor}{\mkbibbold\bgroup}{}{}
\xapptobibmacro{bbx:editor}{\egroup}{}{}

% Grafiken ---------------------------------------------------------------------
% Einbinden von JPG-Grafiken ermöglichen
\usepackage[dvips,final]{graphicx}
\graphicspath{{Bilder/}}

% Paket zur Verwendung zusätzlicher Positionsbefehle
\usepackage{float}

% Befehle aus AMSTeX für mathematische Symbole z. B. \boldsymbol \mathbb --------
\usepackage{amsmath,amsfonts}

% Einfache Definition der Zeilenabstände und Seitenränder etc. -----------------
\usepackage{setspace}
\usepackage{geometry}

% zum Umfließen von Bildern ----------------------------------------------------
\usepackage{floatflt}

% Farbdefinitionen
\usepackage{xcolor} 
\definecolor{hellgelb}{rgb}{1,1,0.9}
\definecolor{colKeys}{rgb}{0,0,1}
\definecolor{colIdentifier}{rgb}{0,0,0}
\definecolor{colComments}{rgb}{1,0,0}
\definecolor{colString}{rgb}{0,0.5,0}
\definecolor{whgreen}{RGB}{113,177,41}
\definecolor{dvBlue}{RGB}{0,85,160}

% URL verlinken, lange URLs umbrechen etc. -------------------------------------
\usepackage{url}

% PDF-Optionen -----------------------------------------------------------------
\usepackage[
    bookmarks,
    bookmarksopen=true,
    colorlinks=true,
% diese Farbdefinitionen zeichnen Links im PDF farblich aus
    linkcolor=dvBlue, % einfache interne Verkn�pfungen
    anchorcolor=black,% Ankertext
    citecolor=dvBlue, % Verweise auf Literaturverzeichniseintr�ge im Text
    filecolor=magenta, % Verkn�pfungen, die lokale Dateien �ffnen
    menucolor=dvBlue, % Acrobat-Men�punkte
    urlcolor=dvBlue, 
% diese Farbdefinitionen sollten für den Druck verwendet werden (alles schwarz)
%    linkcolor=black, % einfache interne Verkn�pfungen
%    anchorcolor=black, % Ankertext
%    citecolor=black, % Verweise auf Literaturverzeichniseintr�ge im Text
%    filecolor=black, % Verkn�pfungen, die lokale Dateien �ffnen
%    menucolor=black, % Acrobat-Men�punkte
%    urlcolor=black, 
    %backref, % Inkompatibel mit BibLateX
    plainpages=false, % zur korrekten Erstellung der Bookmarks
    pdfpagelabels, % zur korrekten Erstellung der Bookmarks
    hypertexnames=false, % zur korrekten Erstellung der Bookmarks
    linktocpage % Seitenzahlen anstatt Text im Inhaltsverzeichnis verlinken
]{hyperref}

% fortlaufendes Durchnummerieren der Fußnoten ----------------------------------
\usepackage{chngcntr}
%\counterwithout{footnote}{chapter}

% Formatierung von Listen ändern -----------------------------------------------
\usepackage{paralist} % itemize, enumerate

% bei der Definition eigener Befehle benötigt
\usepackage{ifthen} % Vielleicht nicht nötig

% definiert u.a. die Befehle \ und \listoftodos
\usepackage{todonotes}
\reversemarginpar

% sorgt dafür, dass Leerzeichen hinter parameterlosen Makros nicht als Makroendezeichen interpretiert werden
\usepackage{xspace}

\usepackage{tabularx} % Tabellenspalten mit variabler Breite
\usepackage{wrapfig}  % Schriftumflossene Bilder


% Seitenstil

% Zeilenabstand 1,5 Zeilen
\onehalfspacing

% Seitenränder
% bottom is 20 + X, weil Footer nicht berücksichtigt wird
\geometry{paper=a4paper,left=30mm,right=20mm,top=20mm, bottom=38mm, footskip=8mm}

% Kopf- und Fußzeilen ----------------------------------------------------------
% Kopf- und Fußzeile auch auf Kapitelanfangsseiten
%\renewcommand*{\chapterpagestyle}{scrheadings} 
% Schriftform der Kopfzeile
%\renewcommand{\headfont}{\normalfont}

% Kopfzeile
\ihead{\headmark} % links
\chead{}
%\ohead{\includegraphics[scale=0.1]{Bilder/DJLogo.png}} % rechts
\setlength{\headheight}{20mm} % Höhe der Kopfzeile
% Kopfzeile über den Text hinaus verbreitern
%\setheadwidth[0pt]{textwithmarginpar} 
\setheadsepline[text]{0.4pt} % Trennlinie unter Kopfzeile

% Fußzeile
%\ifoot{} % links
\cfoot{} % mitte
\ofoot{\pagemark} % rechts
%\setfootsepline[text]{0.4pt} % Trennlinie unter Kopfzeile



% Schusterjungen und Hurenkinder vermeiden
%\clubpenalty = 10000
%\widowpenalty = 10000 
%\displaywidowpenalty = 10000

% Verringert den Abstand über den Überschriften
%\renewcommand*{\chapterheadstartvskip}{\vspace*{-\topskip}}

% Irgendwas mit Silbentrennung in MonoSpace (texttt)
\newcommand{\origttfamily}{}% sollte noch nicht definiert sein!
\let\origttfamily=\ttfamily % alte Definition von \ttfamily sichern
\renewcommand{\ttfamily}{\origttfamily \hyphenchar\font=`\-}


% Eigene Befehle und typographische Auszeichnungen für diese

\newcommand{\bs}{$\backslash$}

% einige Befehle zum Zitieren --------------------------------------------------
\newcommand{\Zitat}[2][\
empty]{\ifthenelse{\equal{#1}{\empty}}{\citep{#2}}{\citep[#1]{#2}}}

% zum Ausgeben von Autoren
\newcommand{\AutorName}[1]{\textsc{#1}}
\newcommand{\Autor}[1]{\AutorName{\citeauthor{#1}}}

% Produktnamen
\newcommand{\produkt}[1]{\textbf{#1}}

\newcommand{\code}[1]{\texttt{#1}}

% zum Einbinden von Programmcode -----------------------------------------------
\usepackage{listings}
\usepackage{xcolor}
\usepackage{textcomp}
\lstset{
    float=hbp,
    basicstyle=\ttfamily\color{black}\small,
    identifierstyle=\color{colIdentifier},
    keywordstyle=\color{colKeys},
    stringstyle=\color{colString},
    commentstyle=\color{colComments},
    %columns=flexible,
    tabsize=2,
    frame=single,
    extendedchars=true,
    showspaces=false,
    showstringspaces=false,
    numbers=left,
    numberstyle=\ttfamily\small,
    numbersep=5pt,
    breaklines=true,
    backgroundcolor=\color{hellgelb},
    breakautoindent=true
}

\lstset{literate=
  {á}{{\'a}}1 {é}{{\'e}}1 {í}{{\'i}}1 {ó}{{\'o}}1 {ú}{{\'u}}1
  {Á}{{\'A}}1 {É}{{\'E}}1 {Í}{{\'I}}1 {Ó}{{\'O}}1 {Ú}{{\'U}}1
  {à}{{\`a}}1 {è}{{\'e}}1 {ì}{{\`i}}1 {ò}{{\`o}}1 {ù}{{\`u}}1
  {À}{{\`A}}1 {È}{{\'E}}1 {Ì}{{\`I}}1 {Ò}{{\`O}}1 {Ù}{{\`U}}1
  {ä}{{\"a}}1 {ë}{{\"e}}1 {ï}{{\"i}}1 {ö}{{\"o}}1 {ü}{{\"u}}1
  {Ä}{{\"A}}1 {Ë}{{\"E}}1 {Ï}{{\"I}}1 {Ö}{{\"O}}1 {Ü}{{\"U}}1
  {â}{{\^a}}1 {ê}{{\^e}}1 {î}{{\^i}}1 {ô}{{\^o}}1 {û}{{\^u}}1
  {Â}{{\^A}}1 {Ê}{{\^E}}1 {Î}{{\^I}}1 {Ô}{{\^O}}1 {Û}{{\^U}}1
  {œ}{{\oe}}1 {Œ}{{\OE}}1 {æ}{{\ae}}1 {Æ}{{\AE}}1 {ß}{{\ss}}1 {Δ}{{$\Delta$}}1
  {ç}{{\c c}}1 {Ç}{{\c C}}1 {ø}{{\o}}1 {å}{{\r a}}1 {Å}{{\r A}}1
  {€}{{\EUR}}1 {£}{{\pounds}}1
}

\usepackage{microtype}

\renewcommand*{\dictumwidth}{.41\textwidth}
\renewcommand*{\dictumrule}{}
\renewcommand*{\dictumauthorformat}[1]{--- #1}
\setkomafont{dictumauthor}{%
\scshape
}

\definecolor{bluekeywords}{rgb}{0,0,1}
\definecolor{greencomments}{rgb}{0,0.5,0}
\definecolor{redstrings}{rgb}{0.64,0.08,0.08}
\definecolor{xmlcomments}{rgb}{0.5,0.5,0.5}
\definecolor{types}{rgb}{0.17,0.57,0.68}

\lstdefinestyle{csharp}
{
    language=[Sharp]C,
    captionpos=b,
    %numbers=left, %Nummerierung
    %numberstyle=\tiny, % kleine Zeilennummern
    showspaces=false,
    showtabs=false,
    breaklines=true,
    showstringspaces=false,
    breakatwhitespace=true,
    escapeinside={(*@}{@*)},
    commentstyle=\color{greencomments},
    morekeywords={partial, var, value, get, set},
    keywordstyle=\color{bluekeywords},
    stringstyle=\color{redstrings},
    basicstyle=\ttfamily\small,
}

\lstdefinestyle{hive}
{
    language=SQL,
    captionpos=b,
    %numbers=left, %Nummerierung
    %numberstyle=\tiny, % kleine Zeilennummern
    showspaces=false,
    showtabs=false,
    breaklines=true,
    showstringspaces=false,
    breakatwhitespace=true,
    escapeinside={(*@}{@*)},
    commentstyle=\color{greencomments},
    morekeywords={bigint, double, decimal, char, string, row, format, lines, location, DELIMITED, FIELDS, TERMINATED, STORED, TEXTFILE,
    	OVERWRITE, DIRECTORY},
    keywordstyle=\color{bluekeywords},
    stringstyle=\color{redstrings},
    basicstyle=\ttfamily\small,
}

\lstdefinestyle{xml}
{
    language=XML,
    captionpos=b,
    %numbers=left, %Nummerierung
    %numberstyle=\tiny, % kleine Zeilennummern
    showspaces=false,
    showtabs=false,
    breaklines=true,
    showstringspaces=false,
    breakatwhitespace=true,
    escapeinside={(*@}{@*)},
    commentstyle=\color{greencomments},
    morekeywords={encoding},
    keywordstyle=\color{bluekeywords},
    stringstyle=\color{redstrings},
    basicstyle=\ttfamily\small,
}

\colorlet{punct}{red!60!black}
\definecolor{background}{HTML}{EEEEEE}
\definecolor{delim}{RGB}{20,105,176}
\colorlet{numb}{magenta!60!black}

\lstdefinelanguage{json}{
    basicstyle=\normalfont\ttfamily,
    numbers=left,
    numberstyle=\scriptsize,
    stepnumber=1,
    numbersep=8pt,
    showstringspaces=false,
    breaklines=true,
    frame=single,
    literate=
     *{0}{{{\color{numb}0}}}{1}
      {1}{{{\color{numb}1}}}{1}
      {2}{{{\color{numb}2}}}{1}
      {3}{{{\color{numb}3}}}{1}
      {4}{{{\color{numb}4}}}{1}
      {5}{{{\color{numb}5}}}{1}
      {6}{{{\color{numb}6}}}{1}
      {7}{{{\color{numb}7}}}{1}
      {8}{{{\color{numb}8}}}{1}
      {9}{{{\color{numb}9}}}{1}
      {:}{{{\color{punct}{:}}}}{1}
      {,}{{{\color{punct}{,}}}}{1}
      {\{}{{{\color{delim}{\{}}}}{1}
      {\}}{{{\color{delim}{\}}}}}{1}
      {[}{{{\color{delim}{[}}}}{1}
      {]}{{{\color{delim}{]}}}}{1},
}

%\usepackage{caption} 
%\captionsetup[table]{skip=100pt}
% Bindet die Literaturdaten ein!
\bibliography{Bibliographie}
\begin{document}

%Startstruktur
\setcounter{secnumdepth}{3}
\setcounter{tocdepth}{2}

\pagestyle{empty}
\thispagestyle{plain}
\begin{titlepage}

\begin{center}

\includegraphics[scale=0.7]{Deckblatt/WH_Logo.jpg}

\vspace{2cm}

\Huge{\textbf{Autovermietung (DSL)}}\\[1.5ex]
\Large{\textbf{Projektdokumentation}}
\rule{\textwidth}{0.4pt}\\[3.0ex]

\large{\textbf{im Masterstudiengang Verteilte Systeme}}\\[3.0ex]

\normalsize
\begin{tabular}{ll}\\
	vorgelegt von: 
	& \quad Daniel Hardes, Fabian Paus \\[1.2ex]
	& \quad \\[1.2ex]
	Modul:  & \quad Business Process Management (BPM) \\[1.2ex]
	Gutachter:  & \quad Prof. Dr. Jürgen Priemer \\[1.2ex]
	Abgabetermin:  & \quad 2. Februar 2016\\[1.2ex]
\end{tabular}

\end{center}

\end{titlepage}


\tableofcontents
\setcounter{page}{1}

\pagestyle{scrheadings}

\newpage

\section{Regeln und Workflows}

Hier wandeln wir die Aussagen in Regeln und Workflows.

\section{Flussdiagramm}

Hier das Flussdiagramm der kompletten Preisfindung.

\section{Einordnung der Regeln}

Die Regeln in die Regelklassen einordnen.

\section{Anwendungssituation}

Wir analysieren die Situation am 21. März 2016 in der Autovermietungsstation.
Dazu beschreiben wir zunächst die Ausgangssituation und dann die einzelnen Kundenanfragen.

\subsection{Ausgangssituation}

Am 21. März 2016 stehen zu jeder Fahrzeugklasse zwei Fahrzeuge zur Verfügung:
\begin{description}
	\item[Kleinklasse:] 2
	\item[Kompaktklasse:] 2
	\item[Mittelklasse:] 2
	\item[Oberklasse:] 2
\end{description}

Relevante Feiertage:
\begin{itemize}
	\item Fr 25.03.2016: Karfreitag
	\item Mo 28.03.2016: Ostermontag
	\item So 01.05.2016: 1. Mai / Tag der Arbeit
\end{itemize}

Klärung:
\begin{itemize}
	\item Wir runden mathematisch, weil das einfacher zu implementieren ist.
	\item Nicht spezifizierte Anfrageparameter nehmen einen Standardwert an (Nein, bzw. 0).
	\item Bei nicht angegebenem Alter ist der Kunde älter als 21 Jahre.
	\item Der Filialleiter genehmigt alle Anfragen.
\end{itemize}

\subsection{Kunde A}

Kundenanfrage:
\begin{itemize}
	\item Klasse: Mittel
	\item Dauer: 5 Tage
	\item Automatik: Nein
	\item Sicherheitstraining: Ja
	
	\item Alter: > 21 Jahre
	\item Neukunde: Nein
	\item Alte Reklamation: Nein
	\item Führerscheindauer: 20 Jahre
\end{itemize}

Zwischenergebnis:
\begin{itemize}
	\item Zeitraum: 4 Wochentage, 0 Wochenende, 1 Feiertag
	\item Fahranfängerprüfung: Bestanden, Zuschlag: 0\%
	\item Verfügbarkeitsprüfung: Bestanden, Nachlass: 0\%
	\item Tagespreis: 70€/Wochentag, 52,50€/Feiertag
	\item Basispreis = $70€ * 4 + 52,50€ * 1 = 332,50€$
	\item Rabatt = 5\% * Basispreis = 16,63€
	\item Endpreis = 315,87€
	\item Gesamtpreis = Endpreis = 315,87€
\end{itemize}

Zustand nach Vermietung:
\begin{description}
	\item[Kleinklasse:] 2
	\item[Kompaktklasse:] 2
	\item[Mittelklasse:] 1
	\item[Oberklasse:] 2
\end{description}

\subsection{Kunde B}

Kundenanfrage:
\begin{itemize}
	\item Klasse: Mittel
	\item Dauer: 1 Tage
	\item Automatik: Nein
	
	\item Alter: 20 Jahre
	\item Sicherheitstraining: Nein
	\item Neukunde: Ja
	\item Alte Reklamation: Nein
	\item Führerscheindauer: 3 Jahre
\end{itemize}

Zwischenergebnis:
\begin{itemize}
	\item Zeitraum: 1 Wochentag, 0 Wochenende oder Feiertage
	\item Fahranfängerprüfung: Genehmigung, Bestanden, Zuschlag: 10\%
	\item Verfügbarkeitsprüfung: Bestanden, Nachlass: 0\%
	\item Tagespreis: 70€/Wochentag, 52,50€/Feiertag
	\item Basispreis = $70€ * 1 + 52,50€ * 0 = 70,00€$
	\item Rabatt = 10,00€
	\item Endpreis = 60,00€
	\item Gesamtpreis = 110\% * Endpreis = 66,00€
\end{itemize}

Zustand nach Vermietung:
\begin{description}
	\item[Kleinklasse:] 2
	\item[Kompaktklasse:] 2
	\item[Mittelklasse:] 0
	\item[Oberklasse:] 2
\end{description}

\subsection{Kunde C}

Kundenanfrage:
\begin{itemize}
	\item Klasse: Mittel
	\item Dauer: 7 Tage
	\item Automatik: Nein
	
	\item Alter: 30 Jahre
	\item Sicherheitstraining: Nein
	\item Neukunde: Nein
	\item Alte Reklamation: Nein
	\item Führerscheindauer: 12 Jahre
\end{itemize}

Zwischenergebnis:
\begin{itemize}
	\item Zeitraum: 4 Wochentag, 2 Wochenende oder Feiertage, 1 Gratistag
	\item Fahranfängerprüfung: Bestanden, Zuschlag: 0\%
	\item Verfügbarkeitsprüfung: Genehmigung, Bestanden, Upgrade auf höhere Klasse, Nachlass: 0\%
	\item Tagespreis: 70€/Wochentag, 52,50€/Feiertag
	\item Basispreis = $70€ * 4 + 52,50€ * 2 + 0€ * 1 = 385,00€$
	\item Rabatt = 0,00€
	\item Endpreis = 385,00€
	\item Gesamtpreis = 100\% * Endpreis = 385,00€
\end{itemize}

Zustand nach Vermietung:
\begin{description}
	\item[Kleinklasse:] 2
	\item[Kompaktklasse:] 2
	\item[Mittelklasse:] 0
	\item[Oberklasse:] 1
\end{description}

\subsection{Kunde D}

Kundenanfrage:
\begin{itemize}
	\item Klasse: Mittel
	\item Dauer: 1 Tage
	\item Automatik: Nein
	
	\item Alter: 35 Jahre
	\item Sicherheitstraining: Nein
	\item Neukunde: Ja
	\item Alte Reklamation: Nein
	\item Führerscheindauer: 10 Jahre
\end{itemize}

Zwischenergebnis:
\begin{itemize}
	\item Zeitraum: 1 Wochentag, 0 Wochenende oder Feiertage
	\item Fahranfängerprüfung: Bestanden, Zuschlag: 0\%
	\item Verfügbarkeitsprüfung: Genehmigung, Bestanden, Upgrade auf höhere Klasse, Nachlass: 0\%
	\item Tagespreis: 70€/Wochentag, 52,50€/Feiertag
	\item Basispreis = $70€ * 1 + 52,50€ * 0 = 70,00€$
	\item Rabatt = 10,00€
	\item Endpreis = 60,00€
	\item Gesamtpreis = 100\% * Endpreis = 60,00€
\end{itemize}

Zustand nach Vermietung:
\begin{description}
	\item[Kleinklasse:] 2
	\item[Kompaktklasse:] 2
	\item[Mittelklasse:] 0
	\item[Oberklasse:] 0
\end{description}

\subsection{Kunde E}

Kundenanfrage:
\begin{itemize}
	\item Klasse: Ober
	\item Dauer: 6 Tage
	\item Automatik: Nein
	
	\item Alter: 46 Jahre
	\item Sicherheitstraining: Nein
	\item Neukunde: Ja
	\item Alte Reklamation: Nein
	\item Führerscheindauer: 20 Jahre
\end{itemize}

Zwischenergebnis:
\begin{itemize}
	\item Zeitraum: 4 Wochentag, 2 Wochenende oder Feiertage
	\item Fahranfängerprüfung: Bestanden, Zuschlag: 0\%
	\item Verfügbarkeitsprüfung: Genehmigung, Abgelehnt, Nachlass: 10\%, Klasse = Mittel
	
	% Es gibt noch eine Oberklasse:
	\item Tagespreis: 90€/Wochentag, 67,50€/Feiertag
	\item Basispreis = $90€ * 4 + 67,50€ * 2 = 495,00€$
	\item Rabatt = 10,00€
	\item Endpreis = 485,00€
	\item Gesamtpreis = 100\% * Endpreis = 485,00€
	
	% Es gibt keine Oberklasse aber eine Mittelklasse
	\item Tagespreis: 70€/Wochentag, 52,50€/Feiertag
	\item Basispreis = $70€ * 4 + 52,50€ * 2 = 385,00€$
	\item Rabatt = 10,00€
	\item Endpreis = 375,00€
	\item Gesamtpreis = 90\% * Endpreis = 337,50€
\end{itemize}

Abgelehnt: Keine Änderung
Zustand nach Vermietung:
\begin{description}
	\item[Kleinklasse:] 2
	\item[Kompaktklasse:] 2
	\item[Mittelklasse:] 0
	\item[Oberklasse:] 0
\end{description}

\subsection{Kunde F}

Kundenanfrage:
\begin{itemize}
	\item Klasse: Ober
	\item Dauer: 4 Tage
	\item Automatik: Nein
	
	\item Alter: 32 Jahre
	\item Sicherheitstraining: Ja
	\item Neukunde: Nein
	\item Alte Reklamation: Ja
	\item Führerscheindauer: 10 Jahre
\end{itemize}

Zwischenergebnis:
\begin{itemize}
	\item Zeitraum: 4 Wochentag, 0 Wochenende oder Feiertage
	\item Fahranfängerprüfung: Bestanden, Zuschlag: 0\%
	\item Verfügbarkeitsprüfung: Genehmigung, Abgelehnt, Nachlass: 0\%
	\item Tagespreis: 90€/Wochentag, 67,50€/Feiertag
	\item Basispreis = $90€ * 4 + 67,50€ * 0 = 360,00€$
	\item Rabatt = 20\% * 360€ + 5\% * 360€ = 72€ + 18€ = 90€
	\item Endpreis = 270,00€
	\item Gesamtpreis = 100\% * Endpreis = 270,00€
	
	% Siehe E
\end{itemize}

Abgelehnt: Keine Änderung
Zustand nach Vermietung:
\begin{description}
	\item[Kleinklasse:] 2
	\item[Kompaktklasse:] 2
	\item[Mittelklasse:] 0
	\item[Oberklasse:] 0
\end{description}

\subsection{Kunde G}

Kundenanfrage:
\begin{itemize}
	\item Klasse: Kompakt
	\item Dauer: 1 Tage
	\item Automatik: Nein
	
	\item Alter: > 21 Jahre
	\item Sicherheitstraining: Nein (wg. Frau)
	\item Neukunde: Nein
	\item Alte Reklamation: Nein
	\item Führerscheindauer: mind. 6 Jahre 
	\item Fahranfänger wg. Frau
\end{itemize}

Zwischenergebnis:
\begin{itemize}
	\item Zeitraum: 1 Wochentag, 0 Wochenende oder Feiertage
	\item Fahranfängerprüfung: Bestanden, Zuschlag: 10\%
	\item Verfügbarkeitsprüfung: Bestanden, Nachlass: 0\%
	\item Tagespreis: 50€/Wochentag, 37,50€/Feiertag
	\item Basispreis = $50€ * 1 + 37,50€ * 0 = 50,00€$
	\item Rabatt = 0€
	\item Endpreis = 50,00€
	\item Gesamtpreis = 110\% * Endpreis = 55,00€
\end{itemize}

Zustand nach Vermietung:
\begin{description}
	\item[Kleinklasse:] 2
	\item[Kompaktklasse:] 1
	\item[Mittelklasse:] 0
	\item[Oberklasse:] 0
\end{description}

\section{Prototyp}

Hier den Prototypen beschreiben.

%\appendix
%\include{Inhalt/Anhang}

%\newpage
%\printbibliography

\end{document}